\documentclass[letterpaper]{article} % DO NOT CHANGE THIS
\usepackage[submission]{aaai2027} % DO NOT CHANGE THIS
\usepackage[hyphens]{url} % DO NOT CHANGE THIS
\usepackage{graphicx} % DO NOT CHANGE THIS
\urlstyle{rm} % DO NOT CHANGE THIS
\def\UrlFont{\rm} % DO NOT CHANGE THIS
\usepackage{natbib} % DO NOT CHANGE THIS AND DO NOT ADD OPTIONS
\usepackage{caption} % DO NOT CHANGE THIS AND DO NOT ADD OPTIONS
\usepackage{amsmath}
\usepackage{amssymb}
\usepackage{booktabs}
\frenchspacing % DO NOT CHANGE THIS

\pdfinfo{
/TemplateVersion (2027.1)
}

\setcounter{secnumdepth}{0}

% Working title. Keep Title Case for the AAAI submission.
\title{Reliability-Aware Biological Joint Moment Estimation for Risk-Sensitive Task-Agnostic Exoskeleton Assistance}
\author{Anonymous Submission}
\affiliations{}

\begin{document}
\maketitle

\begin{abstract}
Task-agnostic exoskeleton controllers estimate continuous biological joint moments to assist movements that are difficult to enumerate in advance.
However, the same direct prediction-to-torque pathway can convert sensor dropout, drift, packet loss, or cross-modal delay into harmful assistance, while clean-data accuracy provides no indication of when an individual prediction should be trusted.
We formulate this deployment problem as reliability-aware biological joint moment estimation under concurrent participant, task, and sensor-fault shift.
Our method augments a causal temporal convolutional network with heteroscedastic moment prediction, synthetic fault supervision, clean-input reconstruction, and short-horizon forecasting.
These learned outputs are complemented by causal integrity statistics that target staleness, disrupted cross-modal coherence, and drift.
Using only held-out validation participants and faults, the method normalizes and selects complementary detection signals and converts causal risk into a continuous assistance gate.
We evaluate participant-disjoint in-distribution and held-out-task trials under the configured training-time fault set and three unseen fault generators.
The protocol jointly measures moment accuracy and calibration, fault AUROC, wrong-direction torque, retained aligned assistance, command smoothness, and an actuator-replay proxy for biological burden.
Relative to a split-matched deterministic TCN, the proposed method increases clean moment RMSE by $0.011$ Nm/kg on ID tasks and $0.021$ Nm/kg on held-out tasks, reduces non-clean wrong-direction assistance by $0.002$ absolute on average, and retains $98.1\%$ of clean aligned torque across three seeds.
This framework treats reliability estimation as a validation-calibrated deployment layer complementary to task generalization and data-efficient domain adaptation.
\end{abstract}

\section{Introduction}

Lower-limb exoskeletons are moving beyond controllers designed for a small catalog of steady locomotion modes.
Everyday assistance must remain useful through transitions and through cyclic, transient, and unstructured movements whose boundaries may be ambiguous.
Task-agnostic control addresses this challenge by replacing discrete activity recognition with a continuous physiological control variable: hip and knee biological joint moments estimated from wearable sensors and mapped directly to assistance \citep{molinaro2024task}.
Related work further reduces the device-specific labels needed to learn such estimators by translating large biomechanics datasets into a target sensor domain \citep{scherpereel2025domain}.
These advances expand the behavioral coverage and data efficiency of learned control.

Behavioral generalization, however, is not the same as run-time reliability.
An estimator may perform well on average across held-out participants and tasks yet encounter a corrupted observation within an otherwise familiar trial.
An encoder can stop updating, an inertial measurement unit (IMU) can drift, packets can arrive in bursts, and independently sampled modalities can become misaligned.
Because moment estimates enter the torque pipeline directly, an unnoticed error can attenuate useful assistance or, more seriously, apply torque opposite to the user's biological moment.
Hardware limits can bound command magnitude, but they do not determine whether the predicted direction is appropriate.
A deployable controller therefore needs both a state estimate and evidence about whether that estimate remains trustworthy.

Standard predictive evaluation does not answer this question.
Root-mean-square error (RMSE) aggregates nominal and failure behavior and does not expose rare but control-critical errors.
Heteroscedastic or epistemic uncertainty can flag unfamiliar inputs, but structured sensor faults need not resemble broad distribution shift: a held value may appear locally plausible, a common-mode delay may preserve amplitude, and slow drift may initially lie inside the training range.
Conversely, an unusual but valid movement can be uncertain without indicating sensor failure.
No single score is therefore expected to cover all fault mechanisms, and a high anomaly score alone does not specify how assistance should change.
This leads to our central question: \emph{can a task-agnostic moment estimator detect complementary forms of unreliability and translate them into a safer control response without unnecessarily suppressing nominal assistance?}

We study this question as reliability-aware causal sequence regression.
A shared temporal convolutional network (TCN) predicts hip and knee moment means and data-dependent variances together with a binary fault probability.
Reconstruction and short-horizon forecasting heads probe whether a corrupted observation remains compatible with the learned sensor dynamics.
Synthetic faults provide supervision while retaining the clean moment and sensor targets whenever possible.
Because learned uncertainty can miss simple physical signatures, we additionally measure causal staleness and use cross-modal coherence and drift as offline diagnostic channels.
The resulting signals are normalized on clean validation data, and a compact detector subset is selected using validation faults only.
For actuation, the causal channels are fused into a timestep-level risk score that continuously attenuates the baseline moment-derived torque after a validation-calibrated deadband.

Our evaluation separates three sources of shift that are often conflated: unseen participants, held-out task families, and unseen fault generators.
The last category---burst packet loss, partial channel loss, and randomized delay jitter---is excluded from both model training and detector selection.
Alongside RMSE, likelihood, and uncertainty calibration, we report fault discrimination and control-oriented outcomes: wrong-direction torque, retained aligned assistance, command smoothness, and an offline actuator-replay proxy.
This protocol tests not only whether a fault is detectable, but whether detection produces a favorable safety--utility trade-off.

Our contributions are fourfold:
\begin{itemize}
    \item We formulate reliability-aware task-agnostic exoskeleton control as a joint estimation, fault-discrimination, and safety--utility problem under participant, task, and sensor-fault shift.
    \item We develop a causal multi-head TCN that couples hip--knee moment estimation with heteroscedastic uncertainty, fault supervision, clean-input reconstruction, and short-horizon forecasting.
    \item We introduce a detectability-aware reliability layer that combines learned evidence with mechanism-specific integrity statistics, freezes calibration and signal selection on validation data, and maps causal risk to a continuous assistance gate.
    \item We establish a participant-disjoint, cross-task evaluation with held-out fault generators and control-oriented metrics that reveal failure modes hidden by clean RMSE alone.
\end{itemize}

\section{Related Work}

\subsection{Task-Agnostic Exoskeleton Control}

Conventional exoskeleton controllers typically first infer a discrete locomotion mode or gait phase and then invoke a task-specific torque policy.
This hierarchy is effective for a small set of structured activities but becomes difficult to enumerate and tune for transient or unstructured movement.
Moment-based control instead uses an instantaneous estimate of the user's biological joint moment as a continuous control state.
An initial hip-exoskeleton study showed that a TCN moment estimator could unify assistance across multiple ambulatory conditions \citep{molinaro2024unifies}.
Molinaro et al. subsequently extended this principle to coordinated hip--knee assistance over cyclic, impedance-like, and unstructured activities without explicit task classification \citep{molinaro2024task}.

These works establish the value of continuous physiological-state estimation, and the latter supplies the sensing configuration, targets, and control transformation used here.
They optimize and evaluate estimation accuracy under nominal operation; their sensor ablations quantify sensitivity after removing a modality but do not infer sensor health, calibrate deployment risk, or change assistance in response to a detected failure.
Our objective is therefore not a new task-agnostic control variable, but a reliability layer for deciding how strongly that variable should be trusted online.

\subsection{Data Efficiency and Domain Shift in Wearable Robots}

Generalization across users and activities depends on diverse motion-capture and force-plate labels, whose collection is costly and device specific.
Scherpereel et al. address this training-time bottleneck with a CycleGAN-based translation through a simulated-sensor domain \citep{scherpereel2025domain}.
Their unsupervised and semisupervised models translate open-source biomechanics data into a device's sensor domain, reducing the labeled target data needed for joint-moment estimation.
Our work addresses a different axis of generalization: run-time task shift and corrupted target-device observations after an estimator has been trained.
Domain adaptation changes how training examples are obtained, whereas our monitor evaluates each deployed sequence and attenuates the resulting command.
The approaches are complementary, and applying the proposed reliability layer to a domain-adapted estimator is a natural extension; we do not claim to reimplement the translation model in this paper.

\subsection{Uncertainty, Time-Series Anomalies, and Selective Prediction}

Heteroscedastic regression separates input-dependent observation noise from the predicted mean \citep{kendall2017uncertainties}.
Monte Carlo dropout and deep ensembles provide practical approximations to epistemic uncertainty \citep{gal2016dropout,lakshminarayanan2017ensembles}, but uncertainty estimates can become miscalibrated under dataset shift \citep{ovadia2019trust}.
For multivariate time series, reconstruction error and disrupted inter-variable associations offer additional anomaly evidence \citep{song2023memto,dai2024sarad}.
These observations motivate our use of several fault-mechanism-specific signals rather than treating predictive variance as a universal detector.

Selective prediction learns to abstain when expected error is high and evaluates the resulting risk--coverage trade-off \citep{geifman2019selectivenet}.
An assistive controller cannot simply discard an isolated prediction without specifying the resulting actuation policy.
We instead map calibrated risk to a continuous gain and select that mapping using control-oriented constraints: preserve aligned assistance under nominal inputs while not increasing wrong-direction torque on validation faults.
This connects uncertainty and anomaly detection to a safety--utility decision at every command timestep.

\section{Problem Formulation}

Let $\mathbf{x}_{1:T}\in\mathbb{R}^{C\times T}$ be a causal window of unilateral foot, shank, and thigh IMUs, insole measurements, and hip--knee encoder states.
The target $\mathbf{y}_{t}\in\mathbb{R}^{2}$ contains mass-normalized hip and knee biological moments.
At deployment, observations may be transformed by an unknown corruption $f\in\mathcal{F}$, producing $\widetilde{\mathbf{x}}_{1:T}$.
We seek a predictor $\widehat{\mathbf{y}}_t$ and reliability gate $g_t\in[0,1]$ such that the commanded torque
\begin{equation}
    \boldsymbol{\tau}^{\mathrm{cmd}}_t
    = g_t\,\mathcal{C}(\widehat{\mathbf{y}}_{1:t})
\end{equation}
retains useful assistance on clean inputs while suppressing harmful commands under task shift or sensor faults.
$\mathcal{C}$ denotes the baseline moment-to-torque scaling, delay, filtering, and saturation pipeline.

This formulation separates three goals: moment fidelity, reliability discrimination, and downstream safety--utility trade-off.
Optimizing or reporting only one of them is insufficient.

\section{Method}

\begin{figure*}[t]
\centering
\setlength{\tabcolsep}{5pt}
\begin{tabular}{ccccc}
\fbox{\shortstack{Wearable signals\\$\widetilde{\mathbf{x}}_{1:t}$}}
& $\longrightarrow$ &
\fbox{\shortstack{Causal TCN\\shared representation}}
& $\longrightarrow$ &
\fbox{\shortstack{Moment mean and variance\\fault, reconstruction, forecast}}\\[6pt]
&& $\downarrow$ && $\downarrow$\\[-2pt]
\fbox{\shortstack{Causal integrity statistics\\staleness, coherence, drift}}
& $\longrightarrow$ &
\fbox{\shortstack{Validation calibration\\and signal selection}}
& $\longrightarrow$ &
\fbox{\shortstack{Risk $r_t$ and gate $g_t$\\torque command}}
\end{tabular}
\caption{Method overview. The estimator predicts joint moments and learned reliability signals. Explicit causal statistics cover fault signatures that can remain inconspicuous to regression uncertainty. Validation-only calibration prevents test-set tuning, and the gate attenuates the baseline assistance command when risk exceeds a nominal deadband.}
\label{fig:overview}
\end{figure*}

\subsection{Reliability-Aware Causal TCN}

Let $\mathbf{c},\mathbf{v}\in\mathbb{R}^{C}$ denote the channelwise mean and standard deviation computed from training participants only.
We normalize each observation as
\begin{equation}
    \mathbf{z}_t=(\mathbf{x}_t-\mathbf{c})\oslash
    \max(\mathbf{v},10^{-6}),
\end{equation}
where $\oslash$ is elementwise division.
Following the baseline \citep{molinaro2024task,bai2018tcn}, the shared encoder contains $L=5$ residual temporal blocks.
Block $\ell$ has two causal convolutions with kernel size $k=5$, dilation $2^{\ell-1}$, weight normalization, ReLU activation, and spatial dropout.
Left padding followed by right-side chomp preserves sequence length without exposing future samples.
The receptive field is
\begin{equation}
    H=1+2(k-1)\sum_{\ell=1}^{L}2^{\ell-1}=249
\end{equation}
samples; the first $H-1$ outputs of each window are consequently excluded from training and evaluation.

Five $1\!\times\!1$ convolutional heads map the causal representation $\mathbf{h}_t$ to moment mean $\boldsymbol{\mu}_t\in\mathbb{R}^{2}$, log variance $\mathbf{s}_t\in\mathbb{R}^{2}$, fault logit $a_t$, clean-input reconstruction $\widehat{\mathbf{z}}_t\in\mathbb{R}^{C}$, and $h$-step forecast $\overline{\mathbf{z}}_t\in\mathbb{R}^{C}$.
The main configuration uses a binary fault head; per-fault logits are reserved for an ablation.
We clamp $\mathbf{s}_t$ to $[-8,5]$ and define the diagonal predictive distribution
\begin{equation}
    p(\mathbf{y}_t\mid\widetilde{\mathbf{x}}_{1:t})
    = \mathcal{N}\!\left(\boldsymbol{\mu}_t,
      \operatorname{diag}(\exp\mathbf{s}_t)\right).
\end{equation}

\subsection{Fault-Augmented Multi-Task Learning}

For every training window, we sample a fault $f$ uniformly from the configured set containing clean, missing-insole, encoder-dropout, IMU-bias, packet-loss, stuck-IMU, and sensor-delay conditions.
An operator $\mathcal{A}_f$ produces $\widetilde{\mathbf{x}}=\mathcal{A}_f(\mathbf{x};\boldsymbol{\xi})$, where $\boldsymbol{\xi}$ contains any random packet mask.
Missing modalities are zeroed, IMU bias is a variance-scaled linear ramp, and packet loss uses previous-value hold.
Sensor delay shifts each modality by a fixed multiple of a base delay and pads its prefix with the first observation, producing cross-modal asynchrony rather than a globally coherent time shift.

Let $m_{jt}$ mark a valid moment label and $b_f=\mathbb{I}[f\neq\mathrm{clean}]$.
The masked heteroscedastic loss is
\begin{equation}
\begin{aligned}
\mathcal{L}_{\mathrm{mom}}
&=\frac{1}{2\sum_{j,t}m_{jt}}\sum_{j,t}m_{jt}\,\ell_{jt},\\
\ell_{jt}&=s_{jt}+(y_{jt}-\mu_{jt})^2\exp(-s_{jt}).
\end{aligned}
\end{equation}
Because the moment head receives $\widetilde{\mathbf{x}}$ but retains the clean target $\mathbf{y}$, augmentation also trains the estimator to remain accurate under faults that preserve sufficient information.
The fault head is optimized with timestepwise binary cross-entropy
\begin{equation}
\begin{aligned}
\mathcal{L}_{\mathrm{fault}}=-\frac{1}{|\Omega|}\sum_{t\in\Omega}
\big[&b_f\log\sigma(a_t)\\
&+(1-b_f)\log(1-\sigma(a_t))\big],
\end{aligned}
\end{equation}
where $\Omega$ contains timesteps with at least one valid joint label.

The auxiliary heads learn the uncorrupted normalized sequence rather than reproduce the corrupted input:
\begin{align}
\mathcal{L}_{\mathrm{rec}} & = \operatorname{MSE}_{\Omega}
\left(\widehat{\mathbf{z}}_t(\widetilde{\mathbf{x}}),\mathbf{z}_t(\mathbf{x})\right),\\
\mathcal{L}_{\mathrm{pred}} & = \operatorname{MSE}_{\Omega}
\left(\overline{\mathbf{z}}_t(\widetilde{\mathbf{x}}),\mathbf{z}_{t+h}(\mathbf{x})\right).
\end{align}
The complete objective is
\begin{equation}
    \mathcal{L}=\mathcal{L}_{\mathrm{mom}}+
    \lambda_f\mathcal{L}_{\mathrm{fault}}+
    \lambda_r\mathcal{L}_{\mathrm{rec}}+
    \lambda_p\mathcal{L}_{\mathrm{pred}},
\end{equation}
with $(\lambda_f,\lambda_r,\lambda_p)=(0.2,0.1,0.1)$.
Reconstruction gradients update the shared encoder in the main model.
For forecasting, we stop gradients at $\mathbf{h}_t$ so the head acts as a temporal probe without allowing its auxiliary loss to degrade the moment representation.

\subsection{Complementary Integrity Signals}

Different faults violate different invariants, so we expose eight candidate reliability channels.
The learned channels are fault probability $q_t^{\mathrm{fault}}=\sigma(a_t)$, mean aleatoric scale $q_t^{\mathrm{ale}}=\frac{1}{2}\sum_j\exp(s_{jt}/2)$, reconstruction residual, forecast residual, and optional epistemic scale.
For a modality group $G$, the residual scores are
\begin{align}
q_t^{\mathrm{rec}} &= \max_G\frac{1}{|G|}\sum_{c\in G}
(\widehat z_{ct}-\widetilde z_{ct})^2,\\
q_t^{\mathrm{pred}} &= \max_G\frac{1}{|G|}\sum_{c\in G}
(\overline z_{ct}-\widetilde z_{c,t+h})^2.
\end{align}
Although the heads are supervised toward clean inputs, inference compares them with the available observation, so an integrity violation creates a large disagreement.
Epistemic scale is the joint-averaged standard deviation of moment means across either stochastic-dropout passes or independently trained ensemble members.

Three explicit statistics target temporal signatures that a learned head may miss.
Staleness is a causal moving average of the fraction of channels satisfying $|\widetilde z_{ct}-\widetilde z_{c,t-1}|<\epsilon$.
Coherence differentiates corresponding foot, shank, and thigh IMU axes, computes rolling zero-lag Pearson correlations for each modality pair, and uses one minus the mean absolute correlation.
Drift fits a least-squares slope to every normalized channel over the evaluation window, averages within each modality, and takes the largest group magnitude.
Coherence and drift are reported as diagnostic detector channels; they are not used by the current command-time gate.

\subsection{Validation-Calibrated Fault Detection}

For fault detection, each channel is averaged over valid timesteps to obtain a window score $Q_k$.
We set its reference $\rho_k$ to the 90th percentile on clean validation windows and fuse a subset $S$ by
\begin{equation}
    R^{\mathrm{det}}(S) = \max_{k\in S} Q_k/\max(\rho_k,10^{-6}).
\end{equation}
To avoid reporting a best-of-test detector, we enumerate subsets of at most three available channels using validation faults only.
For validation-fault AUROCs $A_f(S)$, selection maximizes
\begin{equation}
    J_{\mathrm{det}}(S)=\operatorname{mean}_f A_f(S)
    +0.25\min_f A_f(S)-0.002(|S|-1).
\end{equation}
The second term rewards broad fault coverage and the final term favors a compact monitor.
We freeze $S$ and all $\rho_k$ before ID, OOD, and held-out-generator evaluation.
Burst packet loss, independent partial-channel loss, and randomized delay jitter are excluded from both training and detector selection.

\subsection{Utility-Aware Assistance Gate}

Fault-detection AUROC and control risk serve different purposes.
The command-time risk retains the causal timestep channels
$\mathcal{K}_{\mathrm{gate}}=\{\mathrm{fault},\mathrm{ale},\mathrm{rec},\mathrm{pred},\mathrm{epi},\mathrm{stale}\}$ that are available in a given configuration:
\begin{equation}
    r_t=\max_{k\in\mathcal{K}_{\mathrm{gate}}}
    q_t^k/\max(\rho_k,10^{-6}).
\end{equation}
Risk is converted to a continuous gain
\begin{equation}
g_t = \operatorname{clip}_{[0,1]}\!\left(
1-\frac{r_t-(1+d)}{\gamma}\right),
\end{equation}
where $d$ is a deadband and $\gamma$ controls attenuation softness.
Thus $g_t=1$ inside the nominal region $r_t\leq1+d$ and decreases linearly beyond it.

We determine $d$ and $\gamma$ without test data.
Starting from $d_0=0.1$, we set
\begin{equation}
    d=\min\!\left(d_{\max},\max\left[d_0,
    \operatorname{Quantile}_{0.9}(r^{\mathrm{val\text{-}OOD}})-1\right]\right),
\end{equation}
with $d_{\max}=2$, using clean held-out-task trials from validation participants.
We then search $\gamma\in\{0.25,0.5,1,2,4,8\}$.
Feasible values must retain at least 95\% of clean validation aligned torque, respect the clean wrong-direction tolerance, and not increase wrong-direction torque on the designated validation faults.
Among feasible values, the dimensionless objective weights clean retention, relative fault wrong-direction reduction, and fault retention by $1$, $2$, and $0.25$, respectively.

Finally, the baseline control transform $\mathcal{C}$ scales hip and knee moment estimates by $0.20$ and $0.15$, applies the configured joint delays and a 6-Hz low-pass filter, and clips torque to $0.45$ Nm/kg.
The gated command is $\boldsymbol{\tau}^{\mathrm{cmd}}_t=g_t\mathcal{C}(\boldsymbol{\mu}_{1:t})$.
This policy attenuates an unreliable estimate but does not replace independent actuator saturation or hardware safety limits.

\section{Experimental Protocol}

\subsection{Data and Splits}

We use the released exoskeleton sensor and biological-moment data underlying the task-agnostic controller \citep{molinaro2024task,scherpereel2023dataset}.
The default input has 25 unilateral channels and the targets are right hip and knee moments in Nm/kg.
Splits are participant-disjoint.
Four task families---jumping, cutting, weight lifting, and lunges---are excluded from estimator training and evaluated as task OOD.
All preprocessing statistics, risk references, detector subsets, and gate hyperparameters are estimated without test participants.

\subsection{Comparisons and Ablations}

The primary fair baseline is a deterministic, non-augmented TCN trained on exactly the same participant and task split.
The released Nature checkpoint is reported only as a secondary reference because its original training split is not matched to ours.
We compare deterministic training with and without fault augmentation, probabilistic regression, reconstruction, forecasting, Monte Carlo dropout, and a three-seed deep ensemble.
Action-input profiles test whether desired, measured, execution, or interaction torques improve reliability.
We also report leave-one-subject-out evaluation and fault-severity curves.

\subsection{Metrics}

Moment estimation uses RMSE, MAE, $R^2$, Gaussian NLL, uncertainty--error correlation, predictive-interval coverage, and coverage calibration error.
Reliability uses per-channel and fused clean-versus-fault AUROC for both ID and held-out-task data.

Control-oriented replay applies the same moment-to-assistance mapping to ground truth and predictions.
We measure wrong-direction torque ratio, peak wrong torque, retained aligned torque, command jerk, and mean gate value.
A lightweight actuator proxy adds first-order lag and saturation before computing net biological-moment RMS, work proxies, and the fraction of time that assistance increases rather than decreases moment magnitude.
These proxies support relative controller comparison and are not interpreted as measured metabolic benefit.

\subsection{Implementation and Statistical Analysis}

The main model uses five 80-channel TCN blocks, kernel size 5, spatial dropout 0.15, windows of 768 samples with stride 384, and the first 248 samples masked to remove receptive-field history effects.
Models are trained for 20 epochs with Adam at learning rate $10^{-3}$ using seeds 7, 13, and 23.
The principal analysis aggregates mean and standard deviation across seeds; paired seed-level differences compare the proposed model with the split-matched baseline.
This draft reports descriptive seed-level statistics; formal confidence intervals and multiplicity correction are deferred until the final submission freeze.

\section{Results}

We organize the results around four questions: whether reliability training preserves clean moment estimation, whether the frozen monitor detects faults beyond its training generators, whether detection improves the safety--utility trade-off, and which components account for that behavior.
Unless stated otherwise, entries are means and standard deviations over seeds 7, 13, and 23; all detector and gate choices remain fixed from validation data.

\subsection{Moment Accuracy and Calibration}

Table~\ref{tab:accuracy} compares the proposed model with the split-matched deterministic TCN on uncorrupted trials.
On ID tasks, reliability training changes RMSE from $0.1616\pm0.0008$ to $0.1726\pm0.0017$ Nm/kg (paired difference $+0.0110\pm0.0024$), while the corresponding OOD values are $0.2319\pm0.0050$ and $0.2531\pm0.0023$ Nm/kg (paired difference $+0.0213\pm0.0045$).
The probabilistic model obtains coverage ECEs of $0.033\pm0.017$ and $0.028\pm0.010$ on ID and OOD tasks, respectively.
These results show that the reliability-aware model improves fault robustness at a measurable clean-accuracy cost rather than strictly dominating the split-matched deterministic baseline.
The released Nature checkpoint is retained only as contextual reference because its participant and task split is not matched to ours.

\begin{table}[t]
\centering
\scriptsize
\setlength{\tabcolsep}{3.5pt}
\begin{tabular}{llcccc}
\toprule
Method & Split & RMSE & R2 & NLL & Cov.-ECE\\
\midrule
Deterministic & ID  & $0.1616\pm0.0008$ & $0.7980\pm0.0019$ & -- & --\\
Proposed      & ID  & $0.1726\pm0.0017$ & $0.7696\pm0.0045$ & $-1.453\pm0.034$ & $0.033\pm0.017$\\
Deterministic & OOD & $0.2319\pm0.0050$ & $0.7512\pm0.0107$ & -- & --\\
Proposed      & OOD & $0.2531\pm0.0023$ & $0.7036\pm0.0053$ & $-1.063\pm0.022$ & $0.028\pm0.010$\\
\bottomrule
\end{tabular}
\caption{Clean moment estimation and calibration. RMSE is in Nm/kg. ID and OOD denote seen and held-out task families on unseen participants.}
\label{tab:accuracy}
\end{table}

\subsection{Fault Detection under Task Shift}

Table~\ref{tab:detection} reports representative per-fault AUROC values for the frozen fused detector and its candidate channels.
Across the reported training-time fault types, the fused detector reaches macro-AUROC $0.826$ on ID tasks and $0.760$ on held-out tasks.
For the unseen burst, partial-loss, and delay-jitter generators, the corresponding macro-AUROCs are $0.842$ and $0.778$.
The strongest individual signal is coherence for partial channel loss, reaching $1.000$ AUROC on both ID and OOD tasks, whereas IMU bias remains the most difficult fused case at $0.502$ and $0.504$ AUROC.
This per-mechanism breakdown is essential: fusion is useful only if its aggregate gain does not conceal systematic detector failure.

\begin{table*}[t]
\centering
\scriptsize
\setlength{\tabcolsep}{4pt}
\begin{tabular}{llccccccc}
\toprule
Split & Fault & Fused & Logit & Recon. & Forecast & Stale. & Coherence & Drift\\
\midrule
ID  & Insole missing & $0.928\pm0.013$ & $0.478\pm0.054$ & $0.945\pm0.014$ & $0.929\pm0.013$ & $0.859\pm0.001$ & $0.500\pm0.000$ & $0.319\pm0.002$\\
ID  & IMU bias & $0.502\pm0.001$ & $0.534\pm0.003$ & $0.504\pm0.003$ & $0.504\pm0.002$ & $0.502\pm0.000$ & $0.500\pm0.000$ & $0.514\pm0.000$\\
ID  & Burst packet loss (unseen) & $0.743\pm0.009$ & $0.499\pm0.000$ & $0.504\pm0.001$ & $0.515\pm0.002$ & $0.918\pm0.003$ & $0.057\pm0.008$ & $0.503\pm0.001$\\
OOD & Partial channel loss (unseen) & $0.813\pm0.004$ & $0.499\pm0.000$ & $0.501\pm0.001$ & $0.508\pm0.002$ & $0.961\pm0.001$ & $1.000\pm0.000$ & $0.500\pm0.001$\\
OOD & Delay jitter (unseen) & $0.765\pm0.007$ & $0.481\pm0.002$ & $0.550\pm0.006$ & $0.545\pm0.005$ & $0.498\pm0.000$ & $0.964\pm0.004$ & $0.504\pm0.002$\\
\bottomrule
\end{tabular}
\caption{Representative clean-versus-fault AUROC. ``Unseen'' marks generators excluded from model training and detector selection. The full table reports every preregistered fault.}
\label{tab:detection}
\end{table*}

\subsection{Safety--Utility Trade-off}

Detection quality alone does not establish a safer assistance policy.
Table~\ref{tab:safety} therefore compares the proposed estimator before and after applying the frozen gate.
On clean ID and OOD trials, the gate retains $0.982$ and $0.981$ of aligned assistance, with wrong-direction changes of $-0.0001$ and $-0.0030$.
Across the reported training-time fault types, gating changes the wrong-direction ratio by $-0.0003$ on ID and $-0.0034$ on OOD tasks; for unseen generators, the changes are $-0.0002$ and $-0.0036$.
The associated retained-assistance values quantify the price of this reduction rather than treating lower torque as automatically safer.
Peak wrong torque and command jerk follow the same comparison in the complete results, and actuator-replay outcomes are reported only as offline proxies.

\begin{table}[t]
\centering
\scriptsize
\setlength{\tabcolsep}{2.5pt}
\begin{tabular}{lccccc}
\toprule
Scenario & Wrong-U & Wrong-G & Delta & Ret. & Gate\\
\midrule
ID clean & $0.075$ & $0.075$ & $-0.0001$ & $0.982$ & $0.991$\\
OOD clean & $0.077$ & $0.074$ & $-0.0030$ & $0.981$ & $0.973$\\
ID seen faults & $0.092$ & $0.092$ & $-0.0003$ & $0.918$ & $0.964$\\
ID unseen faults & $0.087$ & $0.086$ & $-0.0002$ & $0.962$ & $0.983$\\
OOD seen faults & $0.102$ & $0.098$ & $-0.0034$ & $0.943$ & $0.960$\\
OOD unseen faults & $0.087$ & $0.083$ & $-0.0036$ & $0.976$ & $0.968$\\
\bottomrule
\end{tabular}
\caption{Utility-aware gating trade-off. U and G denote ungated and gated commands; $\Delta=\mathrm{Wrong}_g-\mathrm{Wrong}_u$ (negative is safer), Ret. is the gated-to-ungated aligned-torque ratio, and Gate is the mean gain.}
\label{tab:safety}
\end{table}

\subsection{Ablations and Stress Tests}

Table~\ref{tab:ablation} follows the preregistered attribution sequence rather than removing several interacting components at once.
Because the component sweep was run for one attribution seed, it is treated as diagnostic evidence rather than as the primary three-seed claim.
Fault augmentation changes macro fault AUROC by $+0.067$ and clean OOD RMSE by $+0.029$ Nm/kg in this sweep.
Adding probabilistic prediction, reconstruction, and forecasting produces successive AUROC changes of $-0.083$, $+0.085$, and $-0.020$; their effects on wrong-direction torque are reported in the final column.
These results distinguish improvements due to exposure to corrupted inputs from those due to a particular reliability head, while also showing that not every auxiliary head improves the single-seed aggregate.

\begin{table}[t]
\centering
\scriptsize
\setlength{\tabcolsep}{3pt}
\begin{tabular}{lcccc}
\toprule
Model & RMSE & AUC & Wrong-D & Ret.\\
\midrule
Deterministic & $0.2309$ & $0.628$ & -- & --\\
+ fault augmentation & $0.2598$ & $0.695$ & $-0.0177$ & $0.731$\\
+ probabilistic head & $0.2472$ & $0.612$ & $-0.0152$ & $0.746$\\
+ reconstruction & $0.2613$ & $0.697$ & $-0.0017$ & $0.785$\\
+ forecasting (proposed) & $0.2578$ & $0.677$ & $-0.0023$ & $0.798$\\
\bottomrule
\end{tabular}
\caption{Ordered component attribution. Fault AUC and control metrics are macro-averaged over the fixed fault set; all rows use the same data split and training budget.}
\label{tab:ablation}
\end{table}

Fault-severity sweeps provide a stricter test than a single operating point.
At the largest packet-loss rate considered, fused AUROC is $0.961$ on ID tasks and $0.914$ on OOD tasks; the corresponding retained-assistance ratios are $0.976$ and $0.976$, with wrong-direction changes of $-0.0001$ and $-0.0032$.
At the largest fixed delay, fused AUROC is $0.899$ on ID tasks and $0.788$ on OOD tasks; the gate retains $0.855$ and $0.966$ of aligned assistance and changes wrong-direction torque by $-0.0012$ and $-0.0051$.
The action-input ablation further tests whether desired, measured, execution, or interaction torque improves detection or creates a device-specific shortcut.
Across leave-one-subject-out folds, the principal safety effect is packet-loss shutdown (wrong-direction deltas $-0.097$ ID and $-0.150$ OOD with zero retained assistance), whereas non-packet faults average only $-0.0029$ and $-0.0099$ wrong-direction change.
This indicates that subject-transfer behavior is conservative and mechanism-specific rather than a broad corrective safety effect.

\section{Discussion}

\subsection{Reliability Is Mechanism Dependent}

The proposed framing treats estimator reliability as a control variable rather than only a confidence visualization.
Its central design implication is that missing channels, frozen streams, asynchronous modalities, slow drift, and behavioral OOD violate different properties of a sensor sequence.
The per-fault comparison tests this interpretation more stringently than macro-AUROC: a useful fusion should improve worst-case coverage rather than only mean discrimination.
In our results, reconstruction residuals are most informative for missing-insole faults, staleness dominates packet-loss faults, and coherence contributes most under partial channel loss and delay faults.
The remaining weakness on IMU bias shows that additional channels do not make every corruption identifiable.

This distinction also explains why task OOD and sensor failure should not be collapsed into one binary anomaly label.
A valid novel movement may deserve assistance even when epistemic uncertainty rises, while a familiar movement with a frozen encoder should trigger attenuation immediately.
Calibrating on clean ID and held-out-task validation trials gives the gate an explicit nominal region for both cases.
Holding unseen participants, tasks, and fault generators out of policy selection then tests three different forms of transfer without using the test set to choose a favorable detector.

\subsection{From Detection to Assistance}

AUROC is threshold-independent and does not measure the consequence of a reliability decision.
The gate instead makes the consequence explicit: it trades retained aligned torque against exposure to wrong-direction torque at every timestep.
The observed trade-off---about a $0.002$ absolute mean reduction in wrong-direction assistance across non-clean main-suite rows while retaining $0.981$ of clean aligned assistance and $0.945$ on faulted rows---should therefore be interpreted as an operating point, not as a universal safety threshold.
Applications with different actuator authority or user tolerance can select another point on the validation Pareto curve without retraining the estimator.

Attenuation is deliberately conservative.
The gate can reduce the magnitude of an estimate judged unreliable, but it cannot repair its direction, identify the failed component with certainty, or decide an optimal fallback torque.
Moreover, uniformly low torque could trivially reduce several risk metrics while eliminating benefit.
Reporting clean and faulty retained assistance, mean gate, jerk, and wrong-direction torque together prevents that degenerate policy from appearing successful.
The result on IMU bias is particularly informative because it exposes a detection limitation rather than merely a calibration or gate-response issue: fused AUROC remains near chance and the gate changes wrong-direction torque by at most $0.003$ absolute.

\subsection{Relationship to Data and Task Generalization}

The comparison with the two baseline papers clarifies scope.
Task-agnostic moment estimation expands behavioral coverage \citep{molinaro2024task}; domain adaptation expands access to training data \citep{scherpereel2025domain}; reliability-aware gating aims to prevent silent failures after such estimators are deployed.
The components can therefore be combined rather than viewed as mutually exclusive controllers.
In particular, the present monitor could be attached to a domain-adapted moment estimator, although its calibration would need to be repeated in the translated target domain.
The three directions answer different deployment questions: what continuous state to assist, how to obtain sufficient training data, and when the resulting prediction should influence the actuator.

\section{Limitations}

The evaluated faults are synthetic and cannot represent the complete hardware failure distribution.
The actuator replay assumes fixed recorded kinematics and consequently supports relative safety comparisons, not claims about human adaptation or metabolic benefit.
Moment labels and sensing originate from one exoskeleton platform, and the current implementation does not jointly train the domain-adaptation model from \citet{scherpereel2025domain}.
Finally, a gate can reduce exposure to a detected failure but cannot guarantee safety under an undetectable or adversarial fault; independent low-level torque and hardware limits remain necessary.

\section{Ethical Statement}

Reliability-aware assistance may reduce risk from silent sensing failures, but overconfident claims could encourage premature deployment around human users.
We therefore distinguish offline replay from human-subject evidence, report negative and weak detector results, and avoid presenting the learned gate as a formal safety guarantee.
Released participant data must be handled under the terms and consent conditions of the source dataset.

\section{Conclusion}

We presented a reliability-aware extension of task-agnostic biological joint moment estimation that combines probabilistic sequence modeling, fault supervision, temporal integrity signals, validation-only calibration, and risk-sensitive assistance gating.
The accompanying benchmark evaluates whether a model remains accurate, recognizes when it is unreliable, and converts that recognition into a favorable safety--utility trade-off under participant, task, and fault shift.

% Anonymous submission: acknowledgments intentionally omitted.

\bibliography{references}

\end{document}
